\newpage
\spread{A translator can learn who acted on whom}
\noindent A fluent answer can reverse an event while preserving every participant and object. We trained an optical writer on pairs of sentences in which exchanging two people, deleting a negation or changing a number requires a corresponding change in the answer image.
\platehere{Figure_190_optical_roles}{Agent and patient in two visual translations}{\textbf{a}, The first held English-to-Chinese role pair in the original order. Both continuations shown here preserve the negation as the names exchange places. \textbf{b}, The first held English passive using the verb \emph{invite}, with its reversal. The first writer produces incomplete or incorrect sentences. The writer trained on source paraphrases expresses both events in active Chinese. All four inputs and eight complete outputs were replayed from their original PNGs on CPU and matched the GPU bitmaps.}
\begin{multicols}{2}
\section*{Changing one fact at a time}
The training set contained 7,288 contrast pairs, including 1,302 derived by editing conversational or editorial examples. Training included ordinary bilingual replay. Two writers started at the same checkpoint and saw the same batches. One used the next-glyph log loss alone. At the first differing target glyph the other also paid $0.15\log(1+\exp(1-s_g+s_f))$, where $s_g$ and $s_f$ score the correct image and its rival. Both generated complete answers without supplied target history at evaluation.
\par Holding out unordered pairs of participants or participant--commodity combinations tested recombination within familiar constructions. The starting writer reproduced 5/144 exact answers. Both continuations produced 143/144, including both members of 71/72 directed contrast pairs. Likelihood alone achieved the same gain.
\section*{Where the branch test stops}
Consider two sentences whose reference images agree up to the first different mark. A source-independent rule at that prefix cannot strictly prefer each mark to the other on the corresponding source. Exact free generation also requires reaching that prefix and finishing the suffix and stop. Lean checks the count bounds for these necessary events.
\par On 72 independently worded requests the contrast writer preferred the correct mark to its rival 64 times with the reference prefix supplied. It preferred the correct mark to every option only 40 times and independently reached the prescribed prefix 31 times. No complete output matched the reference bitmap. Three stopped outputs received meaning credit in both grading orders. An alternative wording can preserve the entire proposition and change its first glyph.
\end{multicols}
\noindent{\small Related contrastive translation evaluation in \href{https://aclanthology.org/2021.blackboxnlp-1.5/}{Vamvas and Sennrich (2021)}.}
\newpage
\spread{Recognizing a translation is easier than writing one}
\noindent Does a compact optical encoder learn that two sentences written in different scripts describe the same event? We asked the same system to retrieve a translation from 128 candidates and to generate an answer without seeing any candidates.
\platehere{Figure_191_recognition_and_writing}{Cross-script retrieval improves while free translation remains difficult}{\textbf{a}, Top-one retrieval on 128 held OPUS-100 sentence pairs. The likelihood and aligned encoders have the same architecture. The projected embedding applies the trained alignment head. \textbf{b}, Stopped outputs credited for complete meaning in both grading orders. Sixteen dialogue and sixteen news families each contribute two requests. OPUS contributes 32 families in both directions. Qwen3-8B Q4 receives text.}
\begin{multicols}{2}
\section*{Matching across scripts}
Of the first 626,047 bilingual rows of the official OPUS-100 training release, 240,000 passed the length, ink inventory, script and deduplication checks. Both continuations saw the same 2,208,000 directed presentations including replay. One also learned to distinguish the paired English and Chinese sentences from other members of its batch using cosine similarity at temperature 0.07. The writer predicts a physical bitmap from its learned output inventory. PNG input needs occupied cells in reading order.
\par The projected reader retrieved 109 correct Chinese and 112 correct English sentences. For a rounded gold--rival margin greater than $2\epsilon$, changing every score by at most $\epsilon$ preserves the winning translation. Lean proves this sufficient condition. A CPU replay from 256 PNGs found 107 and 111 correct decisions respectively certified at 0.01 cosine units. This certifies score perturbations; an image degradation requires its own measured score bound.
\section*{Scoring a candidate backwards}
A list of up to six beam candidates and the stopped greedy output permits a second test. Each image candidate is read backwards across languages by the same writer and scored on the original source, adding half its mean reverse log likelihood to the mean forward score. On 64 OPUS requests the aligned writer received two meaning credits with greedy decoding, two with forward selection and four with reverse selection. None of the 433 candidate images equalled its corresponding OPUS reference, even after text normalization. Other valid phrasings were possible.
\par The aligned fit took 213 seconds on one RTX 4090 D and consumed 28.23 kJ of measured device energy. Reranking added 2.97 seconds and 0.47 kJ across 360 requests. These figures exclude antecedent training and host costs. The reference and foil checks, shared judge family and different prior training histories limit comparisons with Qwen.
\end{multicols}
\noindent{\small \href{https://huggingface.co/datasets/Helsinki-NLP/opus-100}{OPUS-100 corpus and split definitions}.}
\newpage
\spread{Changing the wording can repair a role reversal}
\noindent Training on the same target image with several source phrasings tests whether the writer can distinguish a change of voice from a change of event. A human bilingual benchmark tests how far this ability transfers.
\platehere{Figure_192_paraphrase_transfer}{Complete translation under rewording and domain change}{\textbf{a}, Exact stopped Chinese image outputs on new combinations of familiar actors and verbs in two English voices. \textbf{b}, Complete meanings accepted in both orders on ordinary domains and 64 fresh PAWS-X requests comprising both sentences of 16 non-equivalent pairs in both directions. Canonical and varied training used identical targets and update counts. Qwen3-8B Q4 is the text comparator.}
\begin{multicols}{2}
\section*{Supervising the same event twice}
Two continuations of the broader likelihood writer each received 333,448 directed presentations and 11.40 million target positions. Training mixed 5,986 structured targets, 11,695 targets from PAWS-X and conversational, editorial and OPUS replay. In one condition, an alternative source could replace the canonical source. For PAWS-X only a pair labelled equivalent could supply a different source. Chinese training passages in that corpus are machine translations. Human translations are available on its test set. We excluded 13 test and 10 validation records whose Chinese and English equivalence labels differed.
\par Both writers now produced 144/144 exact answers on the original controlled events. The automatic evaluator credited one reworded translation of 72 in each condition. It rejected the varied writer's correct Chinese active rendering of the pictured doctor passive in both orders, and rejected the reversal in one order (see the preceding role-reversal images). The fixed votes were retained. On new PAWS-X sources, both writers received zero credits and the text comparator received 46. The judge accepted all 64 references and accepted two contrasting foils in both orders.
\section*{An independently checkable consequence}
After observing the judge discrepancy, we froze an additional test of English voice. For each held pair of actors and a verb, we exchanged their roles, changed polarity and used active and passive forms. We excluded entire families if a source or target repeated a training or previously assessed sentence under normalization. Of 140 eligible families, 118 survived, yielding 944 requests. Both writers reproduced 471/472 active targets exactly. On passive inputs the canonical writer reproduced 0/472 and the varied writer 422/472. All eight conditions passed for 76/118 families with varied sources and none with canonical sources. This follow-up measures a familiar construction on new complete strings. Several families share the same actors or verb.
\par The first new English PAWS-X source describes Wright moving from New York to Chapel Hill. The generated Chinese mentions New York, introduces the year 1000 and loses the destination. Cross-script retrieval and role-preserving generation on controlled events coexist with this failure on ordinary prose. A subsequent test of progress needs unseen human sentences with matched entities and different relations, freely expressed in the other script.
\end{multicols}
\noindent{\small \href{https://aclanthology.org/D19-1382/}{Yang et al., PAWS-X (2019)}. \href{https://github.com/lachlanchen/ImagizedLanguageModel/tree/main/models/translation_paraphrases}{Checkpoint, PNGs and executable examples}.}
